\documentclass[11pt,reqno]{amsart}

\usepackage[a4paper,margin=1in]{geometry}
\usepackage{amsmath,amsthm,amssymb,mathtools}
\usepackage[colorlinks=true,linkcolor=blue,urlcolor=blue,citecolor=red]{hyperref}
\usepackage{enumitem}
\usepackage{microtype}

\newtheorem{theorem}{Theorem}[section]
\newtheorem{proposition}[theorem]{Proposition}
\newtheorem{lemma}[theorem]{Lemma}
\newtheorem{corollary}[theorem]{Corollary}
\newtheorem{assumption}[theorem]{Assumption}
\newtheorem{definition}[theorem]{Definition}
\newtheorem{remark}[theorem]{Remark}

\newcommand{\C}{\mathbb C}
\newcommand{\R}{\mathbb R}
\newcommand{\N}{\mathbb N}
\newcommand{\PP}{\mathbb P}
\newcommand{\EE}{\mathbb E}
\newcommand{\Tr}{\operatorname{Tr}}
\newcommand{\supp}{\operatorname{supp}}
\newcommand{\dist}{\operatorname{dist}}
\newcommand{\op}{\operatorname{op}}
\newcommand{\Id}{\operatorname{Id}}
\newcommand{\HS}{\mathrm{HS}}
\newcommand{\norm}[1]{\left\lVert #1\right\rVert}
\newcommand{\abs}[1]{\left\lvert #1\right\rvert}
\newcommand{\eps}{\varepsilon}

\title[A moment-Hankel reduction of the PPT threshold]{A moment-Hankel reduction of Aubrun's PPT threshold}
\author{}
\date{}

\begin{document}
\maketitle

\begin{abstract}
We present a moment-Hankel reduction for the constant \(4\) in Aubrun's
PPT threshold for random induced states in the balanced Hilbert--Schmidt
regime.  Let \(G_d\) be a \(d^2\times s(d)\) complex Gaussian matrix,
\(\rho_d=G_dG_d^*/\Tr(G_dG_d^*)\), and \(s(d)/d^2\to\lambda\).  Under two
explicit moment inputs, the reduction yields the qualitative threshold:
\(\rho_d^\Gamma\) is asymptotically positive when \(\lambda>4\), and
asymptotically non-positive when \(\lambda<4\).  Above \(4\), one does not
need the full operator-norm edge theorem merely to prove PPT: it is enough to
show that an exact unnormalised centered trace moment of
\[
  \widehat H_d
  =
  \frac{d^2s(d)}{\norm{G_d}_{\HS}^2}
  \left(\frac{(G_dG_d^*)^\Gamma}{s(d)}\right)
  =
  d^2\rho_d^\Gamma
\]
tends to zero:
\[
  \EE\Tr\left[(\widehat H_d-\Id)^{2p_d}\right]\longrightarrow0
\]
for a growing order \(p_d\), typically \(p_d\asymp\log d\).
When the growing moment input has the form
\[
  \EE\Tr((\widehat H_d-\Id)^{2p_d})
  \le C(\log d)^\alpha d^2 q^{c\log d},
  \qquad 0<q<1,
\]
and \(c\log(1/q)>2\), the same reduction gives an explicit polynomial
failure rate, up to a power of \(\log d\).
Below \(4\), a finite Stieltjes--Hankel certificate detects a negative
eigenvalue: one fixed shifted Hankel determinant has a negative semicircular
limit and concentrates around it.  The sharp non-asymptotic probability
inequalities remain those of Aubrun's stronger edge analysis.
\end{abstract}

\section{Introduction}

Aubrun's theorem says that the partial transpose of a random induced state on
\(\C^d\otimes\C^d\) has a threshold at environment size \(4d^2\).  More
precisely, for a random induced state \(\rho_d\) with environment dimension
\(s(d)\), the event
\[
  \rho_d^\Gamma\ge0
\]
has probability tending to one if \(s(d)/d^2\to\lambda>4\), and probability
tending to zero if \(s(d)/d^2\to\lambda<4\).  Aubrun's original proof is
stronger than this qualitative statement: it gives quantitative estimates on
the probabilities.  The route below is conditional on two moment inputs that
encode the principal analytic work; they are stated explicitly so that the
deterministic certificates can be examined separately.  We keep probability
estimates visible, although the constants and exponents below are not meant
to reproduce the sharp non-asymptotic inequalities of the original theorem.

The point of the reduction is conceptual.  The lower side
\(\lambda<4\) is not proved by locating the smallest eigenvalue directly.
Instead, one finds a finite moment certificate.  If \(\rho_d^\Gamma\ge0\),
then the shifted Hankel matrices built from the moments of
\(d^2\rho_d^\Gamma\) must be positive.  But for every \(\lambda<4\), one can
choose a finite order \(m\) such that the limiting shifted Hankel determinant
is strictly negative.  Concentration of finitely many moments then makes PPT
overwhelmingly unlikely.

The upper side \(\lambda>4\) is complementary.  Let
\[
  H_d=\frac{W_d^\Gamma}{s(d)}.
\]
By Gaussian invariance, \(H_d\) is centered at \(\Id\): \(\EE H_d=\Id\).
For positivity of the random state, however, the exact object is
\[
  \widehat H_d=d^2\rho_d^\Gamma
  =
  \frac{d^2s(d)}{\Tr W_d}H_d.
\]
Thus the normalization by \(\Tr W_d=\norm{G_d}_{\HS}^2\) is not hidden in
the statement.  If \(\widehat H_d\) has a negative eigenvalue, then
\(\widehat H_d-\Id\) has an eigenvalue less than \(-1\).  Hence, for every
\(p\ge1\),
\[
  \widehat H_d\ngeq0
  \quad\Longrightarrow\quad
  \Tr\left((\widehat H_d-\Id)^{2p}\right)\ge1.
\]
Consequently a high centered trace moment bound tending to zero implies PPT
with high probability by Markov's inequality.  This is weaker than proving
operator-norm convergence to the spectral edge, but it is exactly strong
enough for the positivity conclusion.

\section{Model and Statement}

Let \(G_d\in M_{d^2,s(d)}(\C)\) have independent standard complex Gaussian
entries.  Set
\[
  W_d=G_dG_d^*,
  \qquad
  \rho_d=\frac{W_d}{\Tr W_d}.
\]
The matrix \(W_d\), and hence \(\rho_d\), acts on
\(\C^d\otimes\C^d\).  The partial transpose on the second tensor factor is
denoted by \(\Gamma\).

We use the normalized partially transposed matrix
\[
  A_d=d^2\rho_d^\Gamma.
\]
For the upper-side moment estimates we also write
\[
  \widehat H_d=A_d
  =
  \frac{d^2s(d)}{\Tr W_d}
  \left(\frac{W_d^\Gamma}{s(d)}\right).
\]
Its trace is \(d^2\), and its empirical spectral measure has total mass one
after dividing the trace by \(d^2\).  The limiting law is the semicircle law
with center \(1\) and variance \(1/\lambda\), whose support is
\[
  \left[1-\frac{2}{\sqrt\lambda},\,1+\frac{2}{\sqrt\lambda}\right].
\]
This fixed-moment limit is the partial-transpose Wick/shifted-semicircle
input; it is the same moment law that appears in Aubrun's analysis.

\begin{theorem}[Moment-Hankel threshold under explicit inputs]
\label{thm:ppt-threshold}
Let \(s(d)/d^2\to\lambda\in(0,\infty)\).

\begin{enumerate}[label=(\roman*),leftmargin=*]
\item If \(\lambda>4\) and Assumption~\ref{ass:centered-trace-moment}
  holds along the sequence, then
  \[
    \PP\{\rho_d^\Gamma\ngeq0\}\longrightarrow0.
  \]
  More quantitatively, if the explicit logarithmic moment bound
  \eqref{eq:explicit-log-moment-bound} holds with a choice of
  \(p_d=\lfloor c\log d\rfloor\) for which
  \[
    c\log\left(\frac{\lambda}{4+\eta}\right)>2,
  \]
  then, for every
  \[
    0<\delta<c\log\left(\frac{\lambda}{4+\eta}\right)-2,
  \]
  the failure probability is \(O((\log d)^\alpha d^{-\delta})\).
  Equivalently, the scalar endpoint takes a paper-shape bound
  \[
    \EE\Tr((\widehat H_d-\Id)^{2p_d})
    \le C(\log d)^\alpha d^2q^{c\log d},
    \qquad 0<q<1,\quad c\log(1/q)>2,
  \]
  directly to \(\PP\{\rho_d^\Gamma\ngeq0\}\to0\), with the same envelope.
\item If \(\lambda<4\) and Assumption~\ref{ass:moment-concentration} holds,
  then there are constants \(C,c>0\) such that for all
  sufficiently large \(d\),
  \[
    \PP\{\rho_d^\Gamma\ge0\}
    \le C\exp(-c d^2).
  \]
\end{enumerate}
\end{theorem}

\begin{remark}
The upper side stated here is intentionally weaker than Aubrun's full
spectral-edge theorem.  It proves positivity using only a centered trace
moment that tends to zero.  If one supplies Aubrun's strongest
non-asymptotic spectral estimates, the upper-side probability can be upgraded
accordingly.
\end{remark}

\section{Quantitative Inputs}

The proof below is conditional on two explicit inputs.  They are stated
separately to make clear where probability enters, and also to avoid hiding
the hard part of the argument.  The first input is not proved in this article;
it is the upper-side growing-moment theorem that would have to be supplied by a
uniform partial-transpose Wick analysis.  The scalar and deterministic
consequences of that input are proved here.

\begin{assumption}[High-moment input above the threshold]
\label{ass:centered-trace-moment}
Assume \(\lambda>4\) and \(s(d)/d^2\to\lambda\).  With
\[
  H_d=\frac{W_d^\Gamma}{s(d)},
  \qquad
  \widehat H_d
  =
  \frac{d^2s(d)}{\Tr W_d}H_d
  =
  d^2\rho_d^\Gamma,
\]
there is a sequence \(p_d\to\infty\) such that
\[
  \EE\Tr\left((\widehat H_d-\Id)^{2p_d}\right)\longrightarrow0.
\]
\end{assumption}

\begin{remark}
Assumption~\ref{ass:centered-trace-moment} should be read as a target input,
not as a disguised theorem.  The expected source is a growing-order version of
the partial-transpose Wick expansion for the centered matrix \(H_d-\Id\),
together with a radial transfer through
\(\Tr W_d=\norm{G_d}_{\HS}^2\).  Equivalently, one may prove the estimate
directly for \(\widehat H_d-\Id\), which is the exact observable used by the
formal endpoint.  For fixed \(p\), the Wick expansion gives the shifted
semicircle moments, with Catalan leading terms.  To use
\(p=p_d\asymp\log d\), one needs a uniform counting estimate showing that the
non-leading diagrams remain summable in this growing range.  That uniform
counting is the hard spectral-edge content; Aubrun's original proof proves
stronger estimates by carrying out this kind of edge control.

An explicit sufficient version is the following growing-moment estimate.
For some \(\eta>0\) with \(4+\eta<\lambda\), suppose that there are
constants \(C,\alpha,c_0,d_0>0\) such that, for all \(d\ge d_0\) and all
integers \(1\le p\le c_0\log d\),
\begin{equation}
\label{eq:explicit-log-moment-bound}
  \EE\Tr\left((\widehat H_d-\Id)^{2p}\right)
  \le
  C d^2 p^\alpha
  \left(\frac{4+\eta}{s(d)/d^2}\right)^p.
\end{equation}
If \(c_0\log(\lambda/(4+\eta))>2\), then choosing
\(p_d=\lfloor c\log d\rfloor\) with
\[
  \frac{2}{\log(\lambda/(4+\eta))}<c<c_0
\]
gives Assumption~\ref{ass:centered-trace-moment}.  This is the
high-moment combinatorial input.  It is weaker than an operator-norm edge
theorem, but sufficient for PPT.

The estimate \eqref{eq:explicit-log-moment-bound} is deliberately separated
from the rest of the paper.  Fixed-order moment convergence alone does not
imply it: for fixed \(p\), the unnormalised moment still carries the ambient
factor \(d^2\).  The logarithmic order is precisely what can defeat that
factor when \(\lambda>4\).

For formalization it is convenient to separate the random-matrix estimate
from the scalar probability step.  After increasing \(4+\eta\) slightly and
passing to all sufficiently large \(d\), \eqref{eq:explicit-log-moment-bound}
implies a bound of the cleaner form
\begin{equation}
\label{eq:paper-shape-log-bound}
  \EE\Tr\left((\widehat H_d-\Id)^{2p_d}\right)
  \le
  C'(\log d)^\alpha d^2 q^{c\log d},
  \qquad 0<q<1,\quad c\log(1/q)>2.
\end{equation}
The Lean endpoint proves that this scalar envelope alone forces the
non-PPT probability to tend to zero.  The difficult unproved input is the
model-specific growing-moment estimate that supplies such a \(q\).
\end{remark}

For the lower side we need only fixed moments.  Define
\[
  M_{\ell,d}=d^{-2}\Tr(A_d^\ell),
  \qquad \ell\ge0.
\]
Let \(m_\ell(\lambda)\) be the \(\ell\)-th moment of the semicircle law with
center \(1\) and variance \(1/\lambda\), and set
\[
  H_m(\lambda)=\bigl[m_{i+j+1}(\lambda)\bigr]_{0\le i,j\le m}.
\]
The corresponding random Hankel determinant is
\[
  D_{m,d}=\det\bigl[M_{i+j+1,d}\bigr]_{0\le i,j\le m},
  \qquad
  D_m(\lambda)=\det H_m(\lambda).
\]

\begin{assumption}[Fixed moment concentration]
\label{ass:moment-concentration}
For every fixed \(m\), every compact \(J\Subset(0,\infty)\), and every
\(\tau>0\), there are constants \(C,c>0\) and \(d_0\) such that if
\(d\ge d_0\) and \(s/d^2\in J\), then
\[
  \PP\left\{
    \abs{D_{m,d}-D_m(s/d^2)}>\tau
  \right\}
  \le C\exp(-c d^2).
\]
\end{assumption}

\section{The Hankel Threshold Computation}

The finite certificate is governed by a Chebyshev determinant.  Let
\[
  \alpha_m=4\cos^2\left(\frac{\pi}{m+2}\right).
\]
Then \(\alpha_m\uparrow4\).

\begin{proposition}[Chebyshev-Hankel sign]
\label{prop:chebyshev-sign}
For every \(m\ge1\) and every \(\lambda>0\),
\[
  D_m(\lambda)
  =
  (\sqrt\lambda)^{(m+1)^2}
  U_{m+1}\left(\frac{\sqrt\lambda}{2}\right),
\]
where \(U_{m+1}\) is the Chebyshev polynomial of the second kind.  In
particular, \(D_m(\lambda)>0\) for \(\lambda>\alpha_m\).
Moreover, if
\[
  \cos\left(\frac{2\pi}{m+2}\right)
  <
  \frac{\sqrt\lambda}{2}
  <
  \cos\left(\frac{\pi}{m+2}\right),
\]
then \(D_m(\lambda)<0\).  Consequently, for every \(0<\lambda<4\) there is
some \(m\ge1\) such that \(D_m(\lambda)<0\).
\end{proposition}

\begin{proof}
The determinant formula is the closed form for the shifted Hankel determinant
of the centered semicircular moment sequence.  The sign follows from the
trigonometric identity
\[
  U_{m+1}(\cos\theta)=\frac{\sin((m+2)\theta)}{\sin\theta}.
\]
Writing \(\sqrt\lambda/2=\cos\theta\), the first zero occurs when
\(\theta=\pi/(m+2)\), namely at
\(\lambda=4\cos^2(\pi/(m+2))=\alpha_m\).  For
\(\lambda>\alpha_m\), either \(\sqrt\lambda/2\ge1\), where all \(U_n\) are
positive, or \(0<\theta<\pi/(m+2)\), where both sine factors are positive.
If instead
\[
  \cos\left(\frac{2\pi}{m+2}\right)
  <
  \frac{\sqrt\lambda}{2}
  <
  \cos\left(\frac{\pi}{m+2}\right),
\]
then, since cosine is decreasing on \([0,\pi]\),
\[
  \frac{\pi}{m+2}<\theta<\frac{2\pi}{m+2}.
\]
Thus \(\pi<(m+2)\theta<2\pi\), so
\(\sin((m+2)\theta)<0\), while \(\sin\theta>0\).  Hence
\(U_{m+1}(\sqrt\lambda/2)<0\), and therefore \(D_m(\lambda)<0\).

Finally fix \(0<\lambda<4\), and set
\[
  \theta=\arccos\left(\frac{\sqrt\lambda}{2}\right)\in(0,\pi/2).
\]
Choose
\[
  N=\left\lfloor\frac{\pi}{\theta}\right\rfloor+1.
\]
Then \(N\ge3\), and
\[
  \pi<N\theta\le \pi+\theta<\frac{3\pi}{2}.
\]
With \(m=N-2\), one has \(\sin((m+2)\theta)<0\), hence
\(D_m(\lambda)<0\).
\end{proof}

\begin{remark}
The Lean file \texttt{PptFactorization/Threshold.lean} proves exactly this
sign computation through the theorems
\texttt{det\_pos\_above\_threshold} and
\texttt{det\_neg\_below\_threshold}.
\end{remark}

\section{Proof Above Four}

\begin{lemma}[Trace-moment positivity criterion]
\label{lem:trace-moment-positivity}
Let \(H\) be a Hermitian matrix and let \(p\ge1\).  If
\[
  \Tr\left((H-\Id)^{2p}\right)<1,
\]
then \(H\ge0\).
\end{lemma}

\begin{proof}
Let \(y_1,\ldots,y_n\) be the eigenvalues of \(H\).  The eigenvalues of
\((H-\Id)^{2p}\) are \((y_i-1)^{2p}\), hence
\[
  \Tr\left((H-\Id)^{2p}\right)=\sum_{i=1}^n (y_i-1)^{2p}.
\]
If \(H\) had a negative eigenvalue \(y_i<0\), then \(y_i-1<-1\), so
\((y_i-1)^{2p}>1\).  The trace would be at least \(1\), contradicting the
assumption.
\end{proof}

\begin{proof}[Proof of Theorem~\ref{thm:ppt-threshold}(i)]
Assume \(\lambda>4\), and put \(\widehat H_d=d^2\rho_d^\Gamma\).  Since
\(d^2>0\),
\[
  \widehat H_d\ge0
  \quad\Longleftrightarrow\quad
  \rho_d^\Gamma\ge0.
\]
By Lemma~\ref{lem:trace-moment-positivity},
\[
  \{\,\rho_d^\Gamma\ngeq0\,\}
  =
  \{\,\widehat H_d\ngeq0\,\}
  \subseteq
  \left\{
    \Tr\left((\widehat H_d-\Id)^{2p_d}\right)\ge1
  \right\}.
\]
Markov's inequality gives
\[
  \PP\{\rho_d^\Gamma\ngeq0\}
  \le
  \EE\Tr\left((\widehat H_d-\Id)^{2p_d}\right).
\]
Assumption~\ref{ass:centered-trace-moment} makes the right-hand side tend
to zero.

If the explicit bound \eqref{eq:explicit-log-moment-bound} is available,
then for \(p_d=\lfloor c\log d\rfloor\),
\[
  \PP\{\rho_d^\Gamma\ngeq0\}
  \le
  C d^2 p_d^\alpha
  \left(\frac{4+\eta}{s(d)/d^2}\right)^{p_d}.
\]
Since \(s(d)/d^2\to\lambda\) and
\(c\log(\lambda/(4+\eta))>2\), this is
\[
  O\left((\log d)^\alpha d^{-\delta}\right)
\]
for every
\[
  0<\delta<c\log\left(\frac{\lambda}{4+\eta}\right)-2.
\]
In the paper-shape form \eqref{eq:paper-shape-log-bound}, the bound is even
more transparent:
\[
  \PP\{\rho_d^\Gamma\ngeq0\}
  \le C'(\log d)^\alpha d^2q^{c\log d}
  =
  C'(\log d)^\alpha d^{2+c\log q}.
\]
The exponent is negative exactly when \(c\log(1/q)>2\).
\end{proof}

\section{Proof Below Four}

\begin{lemma}[PPT implies shifted Hankel positivity]
\label{lem:ppt-hankel}
If \(\rho_d^\Gamma\ge0\), then for every \(m\ge0\),
\[
  D_{m,d}\ge0.
\]
\end{lemma}

\begin{proof}
If \(\rho_d^\Gamma\ge0\), then \(A_d=d^2\rho_d^\Gamma\) is positive
semidefinite.  Its empirical spectral measure is supported on
\([0,\infty)\).  The matrix
\[
  \bigl[M_{i+j+1,d}\bigr]_{0\le i,j\le m}
\]
is the Gram matrix of the functions
\[
  x^{i}\sqrt{x},
  \qquad 0\le i\le m,
\]
in \(L^2\) of that positive measure.  It is therefore positive
semidefinite, and its determinant is nonnegative.
\end{proof}

\begin{proof}[Proof of Theorem~\ref{thm:ppt-threshold}(ii)]
Assume \(0<\lambda<4\).  By Proposition~\ref{prop:chebyshev-sign}, choose
\(m\ge1\) such that
\[
  D_m(\lambda)<0.
\]
By continuity of \(D_m\), there are \(\gamma>0\) and a compact interval
\(J\Subset(0,\infty)\) containing \(\lambda\), and hence containing
\(s(d)/d^2\) for all large \(d\), such that
\[
  D_m(t)\le -2\gamma
  \qquad (t\in J).
\]
If \(\rho_d^\Gamma\ge0\), Lemma~\ref{lem:ppt-hankel} gives \(D_{m,d}\ge0\).
For large \(d\), since \(s(d)/d^2\in J\), this forces
\[
  \abs{D_{m,d}-D_m(s(d)/d^2)}\ge 2\gamma.
\]
Consequently
\[
  \PP\{\rho_d^\Gamma\ge0\}
  \le
  \PP\left\{
    \abs{D_{m,d}-D_m(s(d)/d^2)}\ge2\gamma
  \right\}.
\]
Assumption~\ref{ass:moment-concentration} gives
\[
  \PP\{\rho_d^\Gamma\ge0\}\le C\exp(-c d^2).
\]
\end{proof}

\section{What This Proves, and What It Does Not Claim}

The proof above recovers the threshold constant \(4\) with probability
estimates of different strength on the two sides:
\[
  \lambda>4
  \quad\Longrightarrow\quad
  \PP\{\rho_d^\Gamma\ge0\}\to1
  \quad\text{under the centered growing-moment input,}
\]
with the explicit rate
\[
  \PP\{\rho_d^\Gamma\ngeq0\}
  \le C(\log d)^\alpha d^2q^{c\log d}
  =
  O((\log d)^\alpha d^{-\delta})
\]
whenever the supplied centered moment bound has
\(0<q<1\), \(c\log(1/q)>2\), and
\(\delta<c\log(1/q)-2\);
and
\[
  \lambda<4
  \quad\Longrightarrow\quad
  \PP\{\rho_d^\Gamma\ge0\}\le Ce^{-c d^2}.
\]
The lower side is a finite-certificate argument.  It is especially robust:
one fixed shifted Stieltjes--Hankel determinant becomes negative in the limit
and concentration does the rest.

The upper side is a growing centered moment argument.  The key observation is
that positivity is easier than full edge convergence: if
\(\Tr((\widehat H_d-\Id)^{2p_d})<1\), then \(d^2\rho_d^\Gamma\) has no
negative eigenvalue.  To obtain the full strength of Aubrun's
non-asymptotic theorem, one still needs the strongest available spectral-edge
probability estimates.  The present article isolates the threshold mechanism
and keeps the hard high-moment input explicit.  In particular,
Assumption~\ref{ass:centered-trace-moment} is not presented as a new proof of
the edge estimate; it is the exact growing-moment statement whose proof
remains to be supplied.  Thus the article is best read as a moment-Hankel
reduction of the constant \(4\), not as a line-by-line replacement for
Aubrun's sharper probability inequalities.

In current formal terms, the scalar and deterministic probability adapters
above the threshold are now closed.  The remaining frontier is not the
Markov step, nor the logarithmic-rate algebra; it is the actual
partial-transpose growing-moment estimate for the exact normalized observable.
The bridge from the canonical finite Gaussian model to the scalar positivity
endpoint is now formalized.

\section{Formalization Notes}

The Lean development contains the algebraic and combinatorial skeleton used
by this article.

\begin{itemize}[leftmargin=*]
\item \texttt{PartialTranspose.lean} defines the concrete
  partial transpose and proves its algebraic identities.
\item \texttt{Threshold.lean} proves the Chebyshev-Hankel sign computation:
  above
  \[
    \alpha_m=4\cos^2(\pi/(m+2)),
  \]
  the determinant is positive, while the negative direction is the sharper
  band statement used in the proof:
  \[
    \cos\left(\frac{2\pi}{m+2}\right)
    <
    \frac{\sqrt\lambda}{2}
    <
    \cos\left(\frac{\pi}{m+2}\right)
    \quad\Longrightarrow\quad
    \det H_m(\lambda)<0.
  \]
  Varying \(m\) gives a negative finite certificate for every
  \(0<\lambda<4\).
\item \texttt{TransferMatrix.lean} records the
  Stieltjes/Hankel certificate: positivity of the partial transpose implies
  positivity of the shifted Hankel matrices.
\item \texttt{AubrunMomentSpine.lean} and
  \texttt{AppendixBAubrunProposition71.lean} formalize the
  exact Wick reduction for the off-diagonal partial transpose and isolate the
  remaining Aubrun counting estimate.
\item \texttt{AubrunAlternative.lean} contains the deterministic bridge used
  above:
  \[
    \texttt{all\_nonneg\_of\_centered\_even\_moment\_lt\_one}.
  \]
  It proves that a finite spectrum is nonnegative once the unnormalised
  centered even moment is \(<1\).  The same file also records the fixed-moment
  bulk adapters: normalized negative-count and negative-mass bounds from
  centered even moments, event-shaped count and mass implications, and the
  scalar Catalan choice
  \[
    \lambda>4,\ \eta>0
    \quad\Longrightarrow\quad
    \exists m,\ \frac{C_m}{\lambda^m}<\eta .
  \]
  The same file also contains the scalar upper-side rate wrappers:
  \[
    \texttt{negative\_event\_measure\_le\_of\_lintegral\_bound\_log\_quadratic\_rpow\_log},
  \]
  which turns a bound of the form
  \(C(\log d)^\alpha d^2q^{c\log d}\) into a finite-\(d\) non-PPT
  probability bound, and
  \[
    \texttt{exists\_log\_order\_constants\_and\_tendsto\_negative\_event\_measure\_zero\_of\_four\_lt},
  \]
  which packages the condition \(\lambda>4\) into constants
  \(0<q<1\) and \(c\) with \(c\log(1/q)>2\).  These theorems audit only to
  the usual Lean foundations; they do not prove the random-matrix
  growing-moment estimate.
\item \texttt{AubrunAlternativeModelBridge.lean} connects the scalar
  positivity endpoint to the exact random induced state normalization.  In
  particular, it contains the Frobenius-mass endpoint
  \[
  \begin{gathered}
    \texttt{tendsto\_rhoGamma\_posSemidef\_measureReal\_one\_of}\\
    \texttt{\_wishartGamma\_frobeniusMass\_traceMoment\_tendsto\_zero},
  \end{gathered}
  \]
  and its varying-space and canonical finite Gaussian versions.  These
  theorems formalize the implication from
  \[
    \EE\Tr\left[
      \left(
        \frac{d^2s(d)}{\norm{G_d}_{\HS}^2}
        \frac{(G_dG_d^*)^\Gamma}{s(d)}
        -\Id
      \right)^{2p_d}
    \right]\longrightarrow0
  \]
  to
  \[
    \PP\{\rho_d^\Gamma\ge0\}\longrightarrow1.
  \]
  The remaining model-specific input is therefore the displayed
  growing-moment convergence itself, plus the harmless eventual
  non-degeneracy of the environment dimension.
\end{itemize}

The formal files therefore support the exact algebraic threshold mechanism.
Above the threshold, the remaining analytic probability input is the growing
centered trace moment estimate.  A full operator-norm edge theorem would be
stronger, but it is not logically necessary for the PPT conclusion.
The fixed-moment convergence/concentration inputs and the growing-moment
combinatorial estimate remain explicit frontiers rather than hidden
formalized facts.

\begin{thebibliography}{9}

\bibitem{Aubrun2012}
G.~Aubrun,
\emph{Partial transposition of random states and non-centered semicircular
distributions},
Random Matrices: Theory and Applications \textbf{1} (2012), no.~2, 1250001.

\bibitem{AubrunSzarek2017}
G.~Aubrun and S.~J. Szarek,
\emph{Alice and Bob Meet Banach: The Interface of Asymptotic Geometric
Analysis and Quantum Information Theory},
Mathematical Surveys and Monographs, vol.~223, American Mathematical Society,
2017.

\bibitem{Ledoux2001}
M.~Ledoux,
\emph{The Concentration of Measure Phenomenon},
Mathematical Surveys and Monographs, vol.~89, American Mathematical Society,
2001.

\end{thebibliography}

\end{document}
